% Ejercicio 3. Obtener la información de las cabeceras de ambos sonidos. 
Usando la función \texttt{str()} se puede inspeccionar las cabeceras de archivos de sonido.
\begin{lstlisting}[caption={Funciones para las cabeceras de los ficheros de sonido}, label={lst:cabeceras_ficheros}]
str(nombre)
str(apellidos)    
\end{lstlisting}

En la siguiente imagen se ve la información obtenida (Figura \ref{fig:cabeceras_fich}).
\begin{figure}[H]
    \centering
    \includegraphics[width=0.95\linewidth]{images/3_cabeceras.png}
    \caption{Cabeceras de los ficheros}
    \label{fig:cabeceras_fich}
\end{figure}

Como ambos sonidos tienen dos canales, el atributo \textit{left} y el atributo \textit{right} valen lo mismo para cada sonido. Además, confirmamos que tienen dos canales porque la variable \textit{stereo} es verdadera.
